\chapter{TOP-D、MirrorOPD 与 vanilla OPD}
\label{chap:topd}

前几章研究了以 KL 为近端项的几何更新。本章研究另一条稳定蒸馏的路径：先在概率空间中构造一个靠近 student 的 teacher，再用冻结行为策略的数据执行多轮内部训练。信任区域策略蒸馏（Trust Region Policy Distillation，TOP-D）\cite{topd26}采用这条路径。它与 MirrorOPD 都区分 outer 和 inner，但控制的是不同的数学量。TOP-D 的算术混合控制负向 log-ratio 尾部；MirrorOPD 的 KL 子问题产生几何混合，并控制 projective log-ratio 误差。

本章从真实训练流程确定哪些量固定，再证明 TOP-D 的方差、算子收敛与有限时域性能下界。每个定理同时说明其对象：token 梯度、理想分布算子或固定 reward MDP。最后将这些条件应用到 MirrorOPD 和 vanilla OPD，区分可继承的结论与需要重新证明的结论。

\section{外部近端 teacher 与内部权重}
\subsection{算术混合及 reward 变换}

记固定 teacher 为 $T$，第 $k$ 个 outer 的行为策略为 $\pi_k$。TOP-D 对每个 prefix $s$ 定义
\begin{equation}
 U_{k,\alpha}(a\mid s)=(1-\alpha)\pi_k(a\mid s)+\alpha T(a\mid s),
 \qquad 0<\alpha<1.
 \label{eq:topd-target}
\end{equation}
该分布已归一化，不需要配分函数。teacher ratio 和蒸馏 reward 为
\begin{equation}
 \rho_k=\frac{T(a\mid s)}{\pi_k(a\mid s)},\qquad
 \widetilde r_{k,\alpha}(s,a)=\log\frac{U_{k,\alpha}(a\mid s)}{\pi_k(a\mid s)}
 =\log(1-\alpha+\alpha\rho_k).
 \label{eq:topd-reward}
\end{equation}
若已有 vanilla log-ratio $r_k=\log\rho_k$，实现只需计算 $\log(1-\alpha+\alpha e^{r_k})$。可以使用 log-sum-exp 避免指数溢出；不需要额外运行第三个 teacher 网络。

对任意 $\rho\geq0$，有
\begin{equation}
 \log(1-\alpha+\alpha\rho)\geq\log(1-\alpha).
 \label{eq:topd-negative-floor}
\end{equation}
当 teacher 给 student 经常采样的 token 极小概率时，vanilla reward 可以趋于负无穷，而 TOP-D 将惩罚限制在有限下界。正向 reward 仍可任意大，因为 $\rho$ 可以任意大。后面的二阶矩定理利用采样概率对大正向 reward 的衰减作用，不假设 reward 双侧有界。

\subsection{一轮完整的 TOP-D 训练}

原论文 Algorithm 1 的动作顺序是：
\begin{enumerate}
 \item outer 开始时复制参数，得到冻结行为策略 $\pi_k=\pi_{\theta_{\mathrm{old}}}$，并采样 prompt batch。
 \item 对每个 prompt 从 $\pi_k$ 生成 $G$ 个 response，记录采样 token 的 old-policy log-prob 和 teacher log-prob。
 \item 计算式\eqref{eq:topd-reward}的 reward，再计算 return 和组内 token advantage。
 \item 在同一批轨迹上执行 $E$ 个内部 epoch。更新 student 参数，保留 old policy、数据、reward、return 和 advantage。
 \item 将最终 student 作为 $\pi_{k+1}$，进入下一 outer，重新复制行为策略并 rollout。
\end{enumerate}

原文采用即时 reward 加长度归一化的未来 reward。对长度 $L_i$ 的第 $i$ 个 response，
\begin{equation}
 \widetilde R_t^i=\widetilde r_t^i+
 \frac{1}{L_i-t}\sum_{j=t+1}^{L_i}\widetilde r_j^i,\qquad t<L_i.
 \label{eq:topd-return}
\end{equation}
终止位置必须定义 $\widetilde R_{L_i}^i=\widetilde r_{L_i}^i$，避免空和除以零。组内统计量由同一 prompt 的所有有效 response token 计算，得到
\[
 \widehat A_t^i=(\widetilde R_t^i-\mu_k)/\sigma_k.
\]
实现还需要规定零方差组的回退和分母稳定项。这些选择影响小 $\alpha$ 的极限行为。

内部目标采用 PPO 形式。令
\[
 w_t^i(\theta)=\frac{\pi_\theta(a_t^i\mid s_t^i)}{\pi_k(a_t^i\mid s_t^i)},
 \qquad N_k=\sum_iL_i,
\]
则 batch 目标为
\begin{equation}
 \widehat L_k(\theta)=\frac1{N_k}\sum_{i,t}
 \min\{w_t^i(\theta)\widehat A_t^i,
 \operatorname{clip}(w_t^i(\theta),1-\epsilon,1+\epsilon)\widehat A_t^i\}.
 \label{eq:topd-ppo}
\end{equation}
这里两种 ratio 必须分开：$\rho_k=T/\pi_k$ 构造 reward，在 batch 内固定；$w(\theta)=\pi_\theta/\pi_k$ 用于内部重要性加权，随参数变化。PPO clip 作用于后者，不是 teacher ratio 的截断。

所以，TOP-D 的外部混合系数 $\alpha$ 是超参数，近端 teacher、reward 与 advantage 在一个 outer 内固定，在下一个 outer 随 $\pi_k$ 刷新。内部 ratio 和 clipping 状态则在参数更新后改变。固定 $\alpha$ 不等于一直训练第一批数据，也不等于所有权重固定。

\section{算术混合与几何 mirror 的方向}

MirrorOPD 的精确 token target 为
\begin{equation}
 G_\eta(p)(a)=\frac{p(a)^{1-\eta}T(a)^\eta}{Z(\eta)},\qquad
 Z(\eta)=\sum_bp(b)^{1-\eta}T(b)^\eta.
 \label{eq:topd-geometric-comparison}
\end{equation}
固定 prefix，TOP-D target 的位移为 $\alpha(T-p)$。几何 target 在零步长处的导数却为
\begin{equation}
 \left.\frac{\dd G_\eta(p)(a)}{\dd\eta}\right|_{\eta=0}
 =p(a)\left[\log\frac{T(a)}{p(a)}-\E_p\log\frac Tp\right].
 \label{eq:topd-geometric-tangent}
\end{equation}
这两个向量一般不同。单独令 $\alpha,\eta$ 很小，不足以让它们具有相同一阶方向。

若 teacher 和 student 本身接近，写 $T(a)=p(a)(1+u_a)$，其中 $\sum_ap(a)u_a=0$。当 $\max_a|u_a|$ 小时，
\[
 \log(T(a)/p(a))=u_a-\tfrac12u_a^2+O(u_a^3).
\]
代入式\eqref{eq:topd-geometric-tangent}，导数等于
\[
 p(a)u_a+O(p(a)\max_b|u_b|^2)=T(a)-p(a)+O(p(a)\max_b|u_b|^2).
\]
两种方向只在 teacher/student 分歧也小的条件下局部接近。容量差距大、ratio 尖锐时不能省掉这个条件。

原论文也讨论了 log-prob 插值：未归一化地设置 $\log\widetilde T=(1-\alpha)\log p+\alpha\log T$，得到 reward $\alpha\log(T/p)$。正确归一化还需要减去 $\log Z$；该常数在单 prefix 上是 action-independent baseline。缩放后的负尾仍无统一下界。几何更新的价值来自 KL 近端结构与 projective 收缩，不能套用 TOP-D 的负向平滑证明。

\section{全概率空间的 token 梯度方差界}
\subsection{定理对象与完整证明}

对随机向量 $g$，定义 $\Var(g)=\E\|g-\E g\|_2^2=\operatorname{tr}\Cov(g)$。TOP-D 的方差定理针对一个 on-policy token 的原始 score-reward 估计量。累积 return、随机组内标准化、内部 importance ratio 和 PPO 分段梯度没有自动包含在该结论中。

\begin{theorem}[TOP-D 原始 token 方差界]
\label{thm:topd-variance}
词表 $\mathcal V$ 有限。给定 prefix，令 $p(a)=\pi_\theta(a\mid s)>0$、$q(a)=T(a\mid s)\geq0$，两者为概率分布。假设相关 prefix、action 和参数上有 $\|\nabla_\theta\log\pi_\theta(a\mid s)\|_2\leq M$。固定 $0<\alpha<1$，从 $a\sim p$ 采样，定义
\[
 \widetilde g=\nabla_\theta\log\pi_\theta(a\mid s)
 \log(1-\alpha+\alpha q(a)/p(a)).
\]
则任意 prefix 分布下都有
\begin{equation}
 \Var(\widetilde g)\leq M^2|\mathcal V|
 \max\{\log^2(1-\alpha),C^\star\alpha\}.
 \label{eq:topd-original-variance}
\end{equation}
其中 $t^\star\in(2,e^2)$ 是 $\log t=2(1-1/t)$ 在 $t>1$ 的唯一非平凡根，$C^\star=(\log t^\star)^2/(t^\star-1)$。不需要 ratio 上界或统一 support floor。
\end{theorem}

\begin{proof}
由方差恒等式，
\[
 \Var(\widetilde g)=\E\|\widetilde g\|_2^2-\|\E\widetilde g\|_2^2
 \leq\E\|\widetilde g\|_2^2.
\]
score 上界将问题化为 reward 二阶矩：
\begin{equation}
 \E[\|\widetilde g\|_2^2\mid s]
 \leq M^2\sum_ap(a)\log^2(1-\alpha+\alpha q(a)/p(a)).
 \label{eq:topd-second-moment-reduction}
\end{equation}
对单个 action 记 $p=p(a),q=q(a)$，令 $h(p,q)=p\log^2(1-\alpha+\alpha q/p)$。

第一种情况是 $q<p$。此时 $1-\alpha\leq1-\alpha+\alpha q/p<1$，因此
\begin{equation}
 h(p,q)\leq p\log^2(1-\alpha)\leq\log^2(1-\alpha).
 \label{eq:topd-negative-case}
\end{equation}
允许 $q=0$，因为 $\alpha<1$ 保证 log 有限。

第二种情况是 $q\geq p$。此时 log 非负，使用 $q\leq1$ 得到
\[
 h(p,q)\leq p\log^2(1-\alpha+\alpha/p)\leq p\log^2(1+\alpha/p).
\]
令 $u=1/p$，右侧为 $f(u)=\log^2(1+\alpha u)/u$。将实际区间 $u\geq1$ 放宽为 $u>0$ 不会破坏上界。直接求导：
\[
 f'(u)=\frac{\log(1+\alpha u)}{u^2}
 \left[\frac{2\alpha u}{1+\alpha u}-\log(1+\alpha u)\right].
\]
令 $t=1+\alpha u>1$，驻点条件为 $w(t)=0$，其中
\[
 w(t)=\log t-2(1-1/t),\qquad w'(t)=(t-2)/t^2.
\]
$w$ 在 $(1,2)$ 严格递减、在 $(2,\infty)$ 严格递增；并且
\[
 w(1)=0,\quad w(2)=\log2-1<0,\quad w(e^2)=2e^{-2}>0.
\]
介值定理与严格单调性给出唯一 $t^\star\in(2,e^2)$。$f'$ 在该根之前为正、之后为负，故 $u^\star=(t^\star-1)/\alpha$ 是全局最大点。代回得
\begin{equation}
 f(u)\leq f(u^\star)=\alpha\frac{(\log t^\star)^2}{t^\star-1}=C^\star\alpha.
 \label{eq:topd-positive-case}
\end{equation}
端点极限 $f(u)\to0$ 当 $u\downarrow0$ 或 $u\to\infty$ 也核实了最大值性质。

两种情况合并，每个 action 的 $h$ 至多为 $B_\alpha=\max\{\log^2(1-\alpha),C^\star\alpha\}$。有限词表求和给出 $\E[\widetilde r_\alpha^2\mid s]\leq|\mathcal V|B_\alpha$。代入式\eqref{eq:topd-second-moment-reduction}，对 prefix 取期望，再使用方差不超过二阶矩，得到结论。
\end{proof}

\subsection{利用归一化得到更紧的界}

\begin{proposition}[不含词表因子的二阶矩界]
\label{prop:topd-sharp-variance}
在上一组假设下，每个 prefix 有
\begin{equation}
 \E[\widetilde r_\alpha^2\mid s]
 \leq\log^2(1-\alpha)+C^\star\alpha\TV(p,q)
 \leq\log^2(1-\alpha)+C^\star\alpha.
 \label{eq:topd-sharper-variance}
\end{equation}
梯度方差至多为右侧乘以 $M^2$ 后对 prefix 取期望。
\end{proposition}
\begin{proof}
负向区域保留式\eqref{eq:topd-negative-case}中的 $p$，得到
\[
 \sum_{q<p}h(p,q)\leq\log^2(1-\alpha)\sum_{q<p}p\leq\log^2(1-\alpha).
\]
正向区域令 $t=1+\alpha(q-p)/p>1$，则
\[
 p\log^2t=\alpha(q-p)\frac{\log^2t}{t-1}\leq C^\star\alpha(q-p).
\]
这里 $\sup_{t>1}\log^2t/(t-1)=C^\star$ 由同一驻点分析得到。$q=p$ 时贡献为零。概率归一化给出
\[
 \sum_{q>p}(q-p)=\tfrac12\sum_a|q(a)-p(a)|=\TV(p,q).
\]
合并两区并用 $\TV\leq1$，即得二阶矩界。梯度控制沿用上一证明。
\end{proof}

该命题是对原论文界的进一步推导。两个界都不要求 ratio 上界。$\alpha\downarrow0$ 时统一界为 $O(\alpha)$；$\alpha\uparrow1$ 时负向项发散。这表示 vanilla 极限失去统一控制机制，不表示每个固定 teacher/student 实例的方差都发散。

\section{小权重的学习速度与标准化}

固定一个 prefix 的分布并令 $\rho=q/p$。对有限 $\rho$，
\begin{equation}
 \log(1+\alpha(\rho-1))
 =\alpha(\rho-1)-\tfrac12\alpha^2(\rho-1)^2+O(\alpha^3).
 \label{eq:topd-small-alpha}
\end{equation}
有限词表且 $p>0$ 时可逐 action 求和，所以未经归一化的期望梯度满足
\begin{align}
 \E_p[\nabla\log p\,\widetilde r_\alpha]
 &=\alpha\sum_a\nabla p(a)(q(a)/p(a)-1)+O(\alpha^2)\\
 &=\alpha\sum_aq(a)\nabla\log p(a)+O(\alpha^2).
 \label{eq:topd-small-alpha-gradient}
\end{align}
最后一步使用 $\sum_a\nabla p(a)=0$。局部主项是 teacher 对数似然方向，幅度按 $\alpha$ 缩小。不调整 optimizer 或预算时，小 $\alpha$ 会减小原始更新信号。

该 Taylor 展开不是对全部概率分布一致的展开：极小 $p(a)$ 可以让 $\rho$ 达到 $1/\alpha$ 或更大。若随机 prefix 没有相应矩条件，不能交换展开与期望。这解释了固定实例常有 $O(\alpha^2)$ 二阶矩，而统一方差界为 $O(\alpha)$。

标准化改变尺度结论。若整个组的 return 乘以同一个 $c>0$，且没有分母稳定项，则
\[
 \frac{cR-c\mu}{c\sigma}=\frac{R-\mu}{\sigma}.
\]
整体 reward 缩小不会自动令标准化 advantage 缩小。TOP-D 的非线性变换还改变 token 相对权重，尤其限制强负向 signal，不能只用纯缩放解释。

有限 batch 的随机 $\mu,\sigma$ 使原始 token 方差定理不能直接用于标准化梯度。纯尺度缩小时，理想标准化精确抵消尺度，不能仅凭 $\sigma$ 变小断言归一化噪声必然增加。若使用稳定分母，有
\[
 \widehat A_c=\frac{c(R-\mu)}{\sqrt{c^2\sigma^2+\varepsilon^2}}.
\]
当 $c\sigma\gg\varepsilon$，尺度近似抵消；当 $c\sigma\ll\varepsilon$，advantage 近似按 $c$ 缩小。组内方差接近零、固定数值误差与稳定项都需单独分析。学习速度应同时比较原始 reward、advantage、梯度和实际策略 KL。

\section{算术近端算子的收敛}

固定 teacher，并定义 $\mathcal A_\alpha(p)=(1-\alpha)p+\alpha T$。选择与 $k$ 无关的 prefix 分布 $\mu$，定义
\[
 d_\mu(p,q)=\E_{s\sim\mu}\|p(\cdot\mid s)-q(\cdot\mid s)\|_1.
\]
同一证明适用于 prompt 上的完整 response 分布，但 token 条件分布的算术混合一般不同于完整 response 分布的算术混合。定理应用前必须固定层级和度量。

\begin{theorem}[TOP-D 算子有限步与渐近误差界]
\label{thm:topd-contraction}
固定 $T$ 与 $0<\alpha\leq1$。实际策略 $p_{k+1}$ 的近端误差为 $e_k=d_\mu(p_{k+1},\mathcal A_\alpha(p_k))$。则
\begin{equation}
 d_\mu(p_K,T)\leq(1-\alpha)^Kd_\mu(p_0,T)
 +\sum_{j=0}^{K-1}(1-\alpha)^{K-1-j}e_j.
 \label{eq:topd-contraction-error}
\end{equation}
若 $\limsup_ke_k\leq e_\infty$，则 $\limsup_Kd_\mu(p_K,T)\leq e_\infty/\alpha$。精确迭代满足 $p_K=T+(1-\alpha)^K(p_0-T)$。
\end{theorem}
\begin{proof}
三角不等式与算术恒等式给出
\begin{align*}
 d_\mu(p_{k+1},T)&\leq e_k+d_\mu(\mathcal A_\alpha(p_k),T)\\
 &=e_k+\E_\mu\|(1-\alpha)(p_k-T)\|_1\\
 &=e_k+(1-\alpha)d_\mu(p_k,T).
\end{align*}
令 $D_k=d_\mu(p_k,T)$、$b=1-\alpha$。逐次代入为
\[
 D_1\leq bD_0+e_0,\quad D_2\leq b^2D_0+be_0+e_1.
\]
若式\eqref{eq:topd-contraction-error}对 $K$ 成立，则
\[
 D_{K+1}\leq bD_K+e_K\leq b^{K+1}D_0+\sum_{j=0}^Kb^{K-j}e_j.
\]
归纳法给出全部有限步界。

当 $0<\alpha<1$，任取 $\delta>0$，极限上确界条件给出 $N$，使 $e_j\leq e_\infty+\delta$ 对所有 $j\geq N$ 成立。拆分误差和：
\[
 \sum_{j<K}b^{K-1-j}e_j
 =\sum_{j<N}b^{K-1-j}e_j+\sum_{j=N}^{K-1}b^{K-1-j}e_j.
\]
第一项只有有限个数，每项随 $K$ 趋零；第二项至多为
\[
 (e_\infty+\delta)\sum_{m=0}^{K-1-N}b^m\leq(e_\infty+\delta)/\alpha.
\]
初始项也趋零。先取极限上确界，再令 $\delta\downarrow0$，得到渐近界。$\alpha=1$ 时单步递推为 $D_{k+1}\leq e_k$，同样成立。精确迭代时 $p_{k+1}-T=b(p_k-T)$，归纳得到最后的分布恒等式。
\end{proof}

该定理属于算术分布算子，并未证明特定 SGD、Adam 或 PPO 能令 $e_k\to0$。优化、参数表示、采样与访问测度误差都包含在 $e_k$ 中。小 $\alpha$ 既降低原始 reward 尾部风险，也使收缩变慢，并把持续误差放大为 $e_\infty/\alpha$。其结构类似 MirrorOPD 的 $\bar\varepsilon/\eta$，但度量不同。

\begin{corollary}[内部 KL 证书]
若 $\E_\mu\KL(p_{k+1}\Vert\mathcal A_\alpha(p_k))\leq\delta_k^2/2$，则 $e_k\leq\delta_k$。
\end{corollary}
\begin{proof}
逐 prefix 用 Pinsker，再对平方根用 Jensen，得到
\[
 e_k\leq\E_\mu\sqrt{2\KL(p_{k+1}\Vert\mathcal A_\alpha(p_k))}
 \leq\sqrt{2\E_\mu\KL(p_{k+1}\Vert\mathcal A_\alpha(p_k))}\leq\delta_k.
\]
\end{proof}
证书要求 KL 与 $e_k$ 使用相同度量。只测当前 rollout KL 时，还需覆盖或重要性校正，才能推出全部目标 prefix 的误差界。共享参数下一个 prefix 的改进也可能伴随另一个 prefix 的退化。

\section{有限时域访问测度与性能下界}
\subsection{固定 reward 与归一化}

考虑相同初始 prompt 分布、相同转移和固定 reward $r(s,a)$ 的 MDP。所有策略在至多 $H$ 个 action 后终止。吸收状态记为 $s_\bot$，reward 为零，所有策略在该状态采用同一虚拟 action。时间纳入状态，以免混淆不同时间的价值函数。

瞬态访问测度和长度为
\begin{equation}
 d^p(s)=\sum_{t=1}^H\Pr_p(s_t=s),\quad
 \ell_p=\sum_{s\ne s_\bot}d^p(s)=\E_pL,\quad \bar d^p=d^p/\ell_p.
 \label{eq:topd-occupancy-def}
\end{equation}
$\bar d^p$ 只在瞬态状态上定义，且要求 $\ell_p>0$。性能记作 $J_r(p)=\E_p\sum_{t=1}^Lr(s_t,a_t)$，避免与 mirror 步长 $\eta$ 混用。固定基准 $p$，其同 reward 价值函数定义优势 $A_r^p=Q_r^p-V_r^p$。对候选 $q$ 定义
\begin{align}
 f_{p,q}(s)&=\E_{a\sim q}A_r^p(s,a),\\
 \zeta_{r,p}(q)&=J_r(p)+\sum_{s\ne s_\bot}d^p(s)f_{p,q}(s),\\
 \xi_{p,q}&=\sup_{s\ne s_\bot}|f_{p,q}(s)|<\infty.
 \label{eq:topd-surrogate-def}
\end{align}
有界有限时域 reward 足以保证 $\xi$ 有限。定理的 reward 不随候选策略变化。

\begin{lemma}[有限时域性能差异恒等式]
\label{lem:topd-performance-difference}
$J_r(q)-J_r(p)=\sum_{s\ne s_\bot}d^q(s)f_{p,q}(s)$。
\end{lemma}
\begin{proof}
由 Bellman 定义，
\[
 A_r^p(s_t,a_t)=r(s_t,a_t)+\E[V_r^p(s_{t+1})\mid s_t,a_t]-V_r^p(s_t).
\]
沿 $q$ 的轨迹取期望并求和，条件全期望和望远镜消去给出
\[
 \E_q\sum_{t=1}^HA_r^p(s_t,a_t)
 =J_r(q)+\E_qV_r^p(s_{H+1})-\E V_r^p(s_1)=J_r(q)-J_r(p).
\]
终止价值为零，初始分布相同。另一方面，按访问测度展开左侧，正好为 $\sum_sd^q(s)f_{p,q}(s)$，得到恒等式。
\end{proof}

\begin{lemma}[访问测度漂移]
\label{lem:topd-occupancy-shift}
任意 $p,q$ 满足
\begin{equation}
 \|d^q-d^p\|_1\leq2H\sum_{s\ne s_\bot}d^p(s)\TV(q(\cdot\mid s),p(\cdot\mid s)).
 \label{eq:topd-occupancy-shift}
\end{equation}
左侧只对瞬态状态求和。
\end{lemma}
\begin{proof}
记 $\nu_t^p,\nu_t^q$ 为包含吸收状态的完整边缘分布。初始分布相同，故 $\nu_1^p=\nu_1^q$。固定 $t>1$，构造混合分布 $P_i$：前 $i$ 个 action 使用 $p$，之后直至第 $t-1$ 个 action 使用 $q$，$i=0,\ldots,t-1$。于是 $P_0(s_t)=\nu_t^q$，$P_{t-1}(s_t)=\nu_t^p$。

相邻 $P_i,P_{i-1}$ 在第 $i$ 个 action 前都使用 $p$，故共同状态边缘为 $\nu_i^p$；第 $i$ 个 action 分别使用 $p,q$，之后的转移与策略又相同。记共同 future Markov 核为 $K_{i,t}(\cdot\mid s,a)$，则
\[
 P_i(s_t\in\cdot)-P_{i-1}(s_t\in\cdot)
 =\sum_s\nu_i^p(s)\sum_a[p(a\mid s)-q(a\mid s)]K_{i,t}(\cdot\mid s,a).
\]
取 $L^1$ 范数，使用核总质量为一，得到
\begin{align*}
 \|P_i(s_t)-P_{i-1}(s_t)\|_1
 &\leq\sum_s\nu_i^p(s)\sum_a|p(a\mid s)-q(a\mid s)|\\
 &=2\E_{\nu_i^p}\TV(p,q).
\end{align*}
该数据处理步骤对随机环境也成立，自回归确定性 prefix 转移只是特例。对混合分布望远镜求和，有
\[
 \|\nu_t^q-\nu_t^p\|_1\leq2\sum_{i=1}^{t-1}\E_{\nu_i^p}\TV(p,q).
\]
限制到瞬态状态不会增大范数。再按时间求和并交换次序：
\begin{align*}
 \|d^q-d^p\|_1&\leq\sum_{t=2}^H\|\nu_t^q-\nu_t^p\|_1\\
 &\leq2\sum_{t=2}^H\sum_{i=1}^{t-1}\E_{\nu_i^p}\TV(p,q)\\
 &=2\sum_{i=1}^{H-1}(H-i)\E_{\nu_i^p}\TV(p,q)\\
 &\leq2H\sum_{i=1}^H\E_{\nu_i^p}\TV(p,q)\\
 &=2H\sum_{s\ne s_\bot}d^p(s)\TV(p(\cdot\mid s),q(\cdot\mid s)).
\end{align*}
最后使用吸收状态上两策略相同，TV 为零。
\end{proof}

\begin{theorem}[TOP-D 使用的有限时域性能下界]
\label{thm:topd-performance-bound}
在上述固定 reward 与有限 horizon 条件下，
\begin{equation}
 J_r(q)\geq\zeta_{r,p}(q)
 -2\xi_{p,q}H\ell_p\E_{\bar d^p}\TV(q(\cdot\mid s),p(\cdot\mid s)).
 \label{eq:topd-performance-bound}
\end{equation}
\end{theorem}
\begin{proof}
性能差异恒等式减去 surrogate 定义，得到
\[
 J_r(q)-\zeta_{r,p}(q)=\sum_{s\ne s_\bot}[d^q(s)-d^p(s)]f_{p,q}(s).
\]
取绝对值并使用 $|f_{p,q}|\leq\xi_{p,q}$，得到
\[
 |J_r(q)-\zeta_{r,p}(q)|\leq\xi_{p,q}\|d^q-d^p\|_1.
\]
代入访问测度漂移引理，再用 $d^p=\ell_p\bar d^p$，得到
\[
 |J_r(q)-\zeta_{r,p}(q)|\leq2\xi_{p,q}H\ell_p\E_{\bar d^p}\TV(q,p).
\]
取绝对误差不等式的负向一侧，即得下界。
\end{proof}

记下界右侧为 $\mathcal M_{r,p}(q)$。由于 $\E_pA_r^p=0$，有 $\mathcal M_{r,p}(p)=J_r(p)$。候选更新若满足 $\mathcal M_{r,p}(q)\geq J_r(p)$，就保证真实固定 reward 回报不降。这是充分条件，有限样本 PPO clipped surrogate 的若干优化步不会自动满足它。PPO clip 也不是全状态 TV 或 KL 硬约束。采样 advantage、长度改造 return 和随机组内 baseline 与真实 $A_r^p$ 的差异需另行控制。

\subsection{固定 reward 回报与蒸馏 KL 的区别}

TOP-D 在一个 outer 冻结 $U_{k,\alpha}$、$\pi_k$ 后，其 reward 为 $\widetilde r_k=\log U_{k,\alpha}-\log\pi_k$，可代入上一性能定理。该 reward 不随内部 student $q$ 变化。以完整序列为例，令 $P_q,P_U,P_{\pi_k}$ 为相应 token 策略诱导的 response 分布。共同 support 且 KL 有限时，链式 log-prob 给出
\begin{align}
 J_{\widetilde r_k}(q)&=\E_{P_q}\log(P_U/P_{\pi_k})\\
 &=-\KL(P_q\Vert P_U)+\KL(P_q\Vert P_{\pi_k}).
 \label{eq:topd-return-kl-identity}
\end{align}
所以增加固定 reward 回报包含接近 teacher 和远离旧策略的两项，不能单凭该回报增加推出逆 KL 减小。若明确优化
\[
 J_{\widetilde r_k}(q)-\KL(P_q\Vert P_{\pi_k})=-\KL(P_q\Vert P_U),
\]
或直接优化目标 KL，才有精确等价性。使用 student 自身的 $\log q$ 构造 reward 又是另一种动态。

原论文将性能下界、内部优化与 KL 证书串成解释框架。严格成立的逻辑链是：实际优化若达到 KL 证书，则算子误差受控；实际策略若提升真实性能下界，则固定 reward 回报单调。达到证书和提升下界仍是条件，empirical clipped objective 本身没有消除它们。

\section{MirrorOPD 可继承的结论}

MirrorOPD 使用 $\log(G_\eta(p)/T)=(1-\eta)\log(p/T)-\log Z$，对 action 取振幅后消去归一化常数。TOP-D 使用 $\mathcal A_\alpha(p)-T=(1-\alpha)(p-T)$。二者都有几何衰减及持续误差放大，但一个是 projective log-ratio 度量，一个是 $L^1$ 度量。度量转换需要 support 条件，不能只令 $\alpha=\eta$ 就互换定理。

\begin{proposition}[有 support floor 时的 log-ratio 方差界]
\label{prop:mirror-conditional-variance}
若 $\pi_k(a\mid s),T(a\mid s)\geq p_{\min}>0$，score 范数不超过 $M$，则任意采样分布下
\[
 g=\nabla_\theta\log\pi_\theta(a\mid s)\log(T(a\mid s)/\pi_k(a\mid s))
\]
满足 $\Var(g)\leq M^2\log^2(1/p_{\min})$。
\end{proposition}
\begin{proof}
两概率都属于 $[p_{\min},1]$，所以 $p_{\min}\leq T/\pi_k\leq1/p_{\min}$。取 log 后有 $|\log(T/\pi_k)|\leq\log(1/p_{\min})$。逐样本梯度平方范数至多为 $M^2\log^2(1/p_{\min})$；取期望并用方差不超过二阶矩即得结论。
\end{proof}

该条件比 TOP-D 的方差假设强。仅缩小 $\eta$ 不会为 $\eta\log(T/\pi_k)$ 创造全局负下界；原始二阶矩若无限，乘以固定正数后仍无限。support floor、reward clipping 和其他尾部条件分别改变假设或算法，应分别说明。

MirrorOPD 在 outer 内固定 $r_k=\log T-\log\pi_k$，取 $p=\pi_k$，可直接应用有限时域性能下界：
\[
 J_{r_k}(q)\geq J_{r_k}(\pi_k)+\ell_k\E_{\bar d^{\pi_k},q}A_{r_k}^{\pi_k}
 -2\xi_kH\ell_k\E_{\bar d^{\pi_k}}\TV(q,\pi_k).
\]
优势必须来自同一固定 reward MDP。即时 token log-ratio 一般不同于含未来项的优势；若用即时 signal，需要限定为 token 局部目标或控制遗漏未来优势的误差。

下一 outer 的 anchor 和 reward 都刷新，故不同 outer 的 $J_{r_k}$ 是不同函数。可以逐 outer 使用固定 reward 下界，但不能将它们串成同一回报的单调链。全局 teacher 匹配由算子收敛定理说明；外部任务性能需另建立 teacher 与任务质量的关系。

\section{vanilla OPD 的极限、反例与结论}
\subsection{三个不同的对象}

vanilla OPD 至少包含三个应分别分析的对象：
\begin{enumerate}
 \item TOP-D 的 $\alpha=1$ reward 极限：$\log(T/\pi_k)$，失去负向平滑边界，内部仍可保留 frozen-old-policy PPO 结构。
 \item MirrorOPD 的 $\eta=1$ 精确子问题：$\E_q\log(T/\pi_k)-\KL(q\Vert\pi_k)=-\KL(q\Vert T)$。全分布、全覆盖且精确求解时，一轮到 teacher。
 \item 神经网络的有限次 SGD/Adam：参数族受限、prefix 采样有限、inner 未必最优、optimizer 有状态，轨迹不等于精确分布算子。
\end{enumerate}
把 vanilla reward 放入训练代码不等于执行精确近端更新。$\eta=1$ 的闭式解确定子问题终点，不能替代 inner 误差和访问覆盖分析。

\begin{proposition}[有限词表、有界 score 仍不足以统一控制 vanilla 方差]
存在二元词表、严格正 teacher 与有界 softmax score，使 vanilla 梯度方差在一族 teacher 上任意大。
\end{proposition}
\begin{proof}
单 prefix 上令 $p_\theta(1)=e^\theta/(1+e^\theta)$，$p_\theta(2)=1/(1+e^\theta)$。在 $\theta=0$，两概率为 $1/2$，score 为 $1/2,-1/2$，范数始终不超过一。对 $L>\log2$ 取
\[
 T_L(1)=e^{-L},\qquad T_L(2)=1-e^{-L}.
\]
teacher 两概率均正。reward 为 $r_1=\log(2e^{-L})$、$r_2=\log(2(1-e^{-L}))$；梯度等概率取 $g_1=r_1/2$、$g_2=-r_2/2$。两点等概率变量方差为 $(g_1-g_2)^2/4$，于是
\[
 \Var(g)=\frac1{16}(r_1+r_2)^2
 =\frac1{16}[\log(4e^{-L}(1-e^{-L}))]^2\longrightarrow\infty.
\]
每个固定 $L$ 的方差仍有限；反例证明全部合法 teacher 上不存在只依赖词表与 score 的一致上界。
\end{proof}

\subsection{teacher 收敛与任务 reward}

匹配 teacher 与任务准确率提升是不同结论。单步二元任务取 $R(1)=1,R(2)=0$。初始 student 对 action 1 的概率为 $0.9$，teacher 为 $0.2$。精确算术与几何迭代都趋于 teacher，任务回报却从 $0.9$ 降到 $0.2$。这是蒸馏目标与外部 reward 的关系，不能由优化定理排除。

若完整序列 reward 满足 $0\leq R(x,y)\leq R_{\max}$，TV 对偶界给出
\[
 |J_R(p)-J_R(T)|\leq R_{\max}\E_x\TV(P_p(\cdot\mid x),P_T(\cdot\mid x)).
\]
证明可将带符号测度 $P_p-P_T$ 分为正负两部分；两部分总质量各为 TV，积分 $[0,R_{\max}]$ 中的函数后，绝对差至多为 $R_{\max}\TV$。再用 Pinsker 和 Jensen，得到
\[
 J_R(p)\geq J_R(T)-R_{\max}\sqrt{\tfrac12\E_x\KL(P_p\Vert P_T)}.
\]
这是相对 teacher 的性能下界，既不保证每步单调，也不保证 teacher 优于初始 student。将 rollout token KL 换成序列 KL 还需链式 KL 与正确的 student 访问测度。

\section{PMD-mean：同一小步长下的不同回归目标}
\label{sec:pmdmean-comparison}

算术平滑、几何近端更新和均值基线回归是三种不同操作。PMD-mean 论文\cite{pmdmean26}说明，把精确 log-partition 换成行为 reward 均值，会改变总体最优解；无限数据不能消除这项算法差异。下面在固定 prefix 或固定 prompt 的有限动作空间中重建该结论。动作可以是 token，也可以是完整 response，但同一推导中必须保持层级不变。

\subsection{精确 target、均值 target 与中心化 residual}

设行为分布 $b(a)>0$，固定信号 $r(a)$，令
\[
 h(a)=r(a)-\E_b r,\qquad \E_bh=0,\qquad \eta>0.
\]
原论文正则化参数为 $\tau$，本书统一用 $\eta=1/\tau$。精确 KL mirror target 和它的 log-ratio 为
\begin{equation}
 q_{\mathrm{part}}(a)=\frac{b(a)e^{\eta h(a)}}{A},\qquad
 A=\E_b e^{\eta h},\qquad
 u_{\mathrm{part}}(a)=\eta h(a)-\log A.
 \label{eq:pmdmean-part-target}
\end{equation}
用精确 $u_{\mathrm{part}}$ 做平方回归，其总体最小 loss 为零，并恢复该 target。

PMD-mean 则删除 $\log A$，拟合 $\eta h$。记候选 $q$ 的 $u(a)=\log(q(a)/b(a))$，得到
\begin{equation}
 L_{\mathrm{mean}}(q)=\tfrac12\E_b[u(a)-\eta h(a)]^2,
 \qquad \E_b e^u=1.
 \label{eq:pmdmean-loss}
\end{equation}
当 $h$ 非常数时，Jensen 严格不等式给出 $A>1$，所以 $u=\eta h$ 不满足归一化。平方回归被迫在无法精确拟合的目标与概率约束间折中。

精确中心化 residual 是另一种 loss：
\begin{equation}
 D_q(a)=\eta h(a)-u(a)-\KL(b\Vert q)
 =\eta h(a)-u(a)+\E_bu,
 \qquad L_{\mathrm{cent}}(q)=\tfrac12\E_bD_q^2.
 \label{eq:pmdmean-centered-loss}
\end{equation}
这里 $\KL(b\Vert q)$ 必须随候选 $q$ 计算，并参与相应优化定义。它不是将精确 log-partition 用固定 reward 均值替代。对 $q_{\mathrm{part}}$，有 $\KL(b\Vert q_{\mathrm{part}})=\log A$，故 $D_{q_{\mathrm{part}}}=0$。

\begin{proposition}[中心化 residual 恢复精确 mirror target]
\label{prop:pmdmean-centered-exact}
有限动作、$b>0$ 时，$L_{\mathrm{cent}}$ 的唯一分布最优解为 $q_{\mathrm{part}}$，最小值为零。在全部自由 logits 参数化中，所有驻点都对应这个分布。
\end{proposition}
\begin{proof}
零 residual 条件是 $u-\E_bu=\eta h$。因此 $u=\eta h+c$，归一化要求 $e^cA=1$，即 $c=-\log A$，得到唯一 $q_{\mathrm{part}}$。平方 loss 非负，所以它是全局最优解。

为证明没有额外驻点，定义零均值坐标 $v=u-\E_bu$。任意 $\E_bv=0$ 的有限向量，都唯一对应归一化分布
\[
 q_v(a)=\frac{b(a)e^{v(a)}}{\E_b e^v}.
\]
反过来每个正分布都对应该零均值坐标，所以这是正单纯形与零均值线性空间之间的光滑一一映射。在此坐标下
\[
 L_{\mathrm{cent}}(q_v)=\tfrac12\E_b(v-\eta h)^2.
\]
对任意零均值扰动 $z$，方向导数为 $\E_b[(v-\eta h)z]$。由于 $v-\eta h$ 也零均值，若所有导数为零，可取 $z=v-\eta h$，得到 $\E_b(v-\eta h)^2=0$。$b>0$ 强制 $v=\eta h$，即唯一 target。全部自由 logits 只增加一个无效常数方向，不增加新的分布驻点。
\end{proof}

该驻点结论依赖完整分布参数化。神经网络参数族、采样 KL 或对 KL 项停止梯度，都需要另分析。特别是 stop-gradient 版本虽然在精确 target 上 residual 仍为零，其全部驻点与真实平方 loss 的导数不再由上面的坐标证明保证。

\subsection{PMD-mean 的 Lambert-$W$ 解}

\begin{theorem}[PMD-mean 总体解及归一化常数]
\label{thm:pmdmean-lambert}
有限动作、$b>0$ 和 $\E_bh=0$ 下，式\eqref{eq:pmdmean-loss}有唯一总体最优分布
\begin{equation}
 q_{\mathrm{mean}}(a)=b(a)\exp\{\eta h(a)-W(\kappa e^{\eta h(a)})\},
 \label{eq:pmdmean-lambert}
\end{equation}
其中 $W$ 为非负实数上的 Lambert-$W$ 主分支，即 $W(z)e^{W(z)}=z$。唯一 $\kappa\geq0$ 由归一化决定。若 $h$ 非常数，$\kappa>0$；若 $h=0$，则 $\kappa=0$ 且 $q_{\mathrm{mean}}=b$。令 $A=\E_be^{\eta h}$、$B=\E_be^{2\eta h}$，则
\begin{equation}
 \frac{A(A-1)}B\leq\kappa\leq\log A.
 \label{eq:pmdmean-kappa-bounds}
\end{equation}
原论文乘子 $\lambda$ 与此处的关系是 $\kappa=\lambda/\tau^2=\lambda\eta^2$。
\end{theorem}
\begin{proof}
在 log-ratio 坐标中，引入归一化乘子：
\[
 \mathcal L(u,\kappa)=\tfrac12\sum_ab(a)[u(a)-\eta h(a)]^2
 +\kappa\left(\sum_ab(a)e^{u(a)}-1\right).
\]
逐 action 求导，并除以正数 $b(a)$，得到
\begin{equation}
 u(a)-\eta h(a)+\kappa e^{u(a)}=0.
 \label{eq:pmdmean-stationarity}
\end{equation}
令 $w_a=\eta h(a)-u(a)$，则 $w_a=\kappa e^{u(a)}$，所以
\[
 w_ae^{w_a}=\kappa e^{\eta h(a)},\qquad
 w_a=W(\kappa e^{\eta h(a)}).
\]
代回得到式\eqref{eq:pmdmean-lambert}。

现在证明存在性、唯一性和全局最优性，而不仅给出形式驻点。定义
\[
 F(\kappa)=\sum_ab(a)\exp\{\eta h(a)-W(\kappa e^{\eta h(a)})\}.
\]
由于 $W$ 严格递增，$F$ 连续且严格递减。$F(0)=A\geq1$；当 $\kappa\to\infty$，每个 summand 趋零，故 $F(\kappa)\to0$。所以存在唯一非负 $\kappa$ 使 $F(\kappa)=1$。$h$ 非常数时 Jensen 给出 $A>1$，根严格正；$h=0$ 时根为零。

固定这一非负乘子，$\mathcal L$ 关于 $u$ 的 Hessian 为对角矩阵
\[
 \frac{\partial^2\mathcal L}{\partial u(a)^2}=b(a)(1+\kappa e^{u(a)})>0.
\]
因此驻点是无约束拉格朗日函数的唯一全局最小点。所有满足归一化的 $u$ 上，乘子项为零，$\mathcal L=L_{\mathrm{mean}}$。驻点本身可行，所以它也是约束问题唯一全局最小点。这一步避免把非凸的指数等式约束误当作凸约束。

为证明乘子界，先设 $\kappa>0$。利用 $e^{-W(z)}=W(z)/z$ 和归一化，有
\begin{equation}
 \kappa=\E_bW(\kappa e^{\eta h}).
 \label{eq:pmdmean-kappa-fixedpoint}
\end{equation}
$W$ 在非负区间递增且凹，Jensen 给出 $\kappa\leq W(\kappa A)$。对两边施加递增函数 $w\mapsto we^w$，得到
\[
 \kappa e^\kappa\leq\kappa A,\qquad\kappa\leq\log A.
\]
另一方面 $W(z)\geq z/(1+z)$。该性质可由 $z=we^w$ 验证：$w\geq z/(1+z)$ 等价于 $1\geq(1-w)e^w$，后者在 $w\geq0$ 上成立，因为右侧导数为 $-we^w\leq0$。

设 $X=e^{\eta h}$，则式\eqref{eq:pmdmean-kappa-fixedpoint}与 Cauchy--Schwarz 给出
\begin{align*}
 \kappa&\geq\E_b\frac{\kappa X}{1+\kappa X}\\
 &\geq\frac{(\E_b\kappa X)^2}{\E_b[\kappa X(1+\kappa X)]}
 =\frac{\kappa A^2}{A+\kappa B}.
\end{align*}
除以 $\kappa$ 并整理，得到 $\kappa\geq A(A-1)/B$。$\kappa=0$ 时 $A=B=1$，两界亦成立。
\end{proof}

\subsection{隐式 KL--$\chi^2$ 正则化的逐步等价证明}

\begin{proposition}[PMD-mean 的混合正则化子问题]
\label{prop:pmdmean-mixed}
上一最优分布也唯一最大化
\begin{equation}
 \E_qh-\frac1\eta\KL(q\Vert b)
 -\frac{\kappa}{2\eta}\chi^2(q\Vert b),\qquad
 \chi^2(q\Vert b)=\sum_ab(a)(q(a)/b(a)-1)^2.
 \label{eq:pmdmean-mixed-objective}
\end{equation}
系数 $\kappa$ 是由回归问题及归一化自适应确定的，不是另一个任意选择的固定超参数。
\end{proposition}
\begin{proof}
在分布变量 $q$ 中，该目标的 Hessian 为对角矩阵，逐项为
\[
 -\frac1{\eta q(a)}-\frac\kappa{\eta b(a)}<0.
\]
所以目标严格凹。带归一化乘子 $\nu$ 的一阶条件为
\[
 h(a)-\frac1\eta(\log(q(a)/b(a))+1)
 -\frac\kappa\eta(q(a)/b(a)-1)+\nu=0.
\]
代入 $u=\log(q/b)$ 并乘以 $\eta$，得
\[
 \eta h(a)-u(a)-1-\kappa(e^{u(a)}-1)+\eta\nu=0.
\]
PMD-mean 的驻点条件\eqref{eq:pmdmean-stationarity}给出 $\eta h-u-\kappa e^u=0$。选择 $\eta\nu=1-\kappa$，上面的混合目标一阶条件逐 action 成立。该分布满足归一化且严格正，严格凹性保证它是唯一最优解。
\end{proof}

这说明 PMD-mean 的保守性来自额外惩罚大概率比率的 $\chi^2$ 项，而不是把精确几何步长换成一个更小常数。精确中心化 residual 没有这项总体偏差；它恢复的是原始 KL 子问题。两者在有限样本中仍会有各自的估计误差。

\subsection{小 $\eta$ 的相似性与总体误差的区别}

有限动作固定 $h$ 时，利用 $\E_bh=0$ 展开得
\[
 A=1+\frac{\eta^2}{2}\E_bh^2+O(\eta^3),\qquad
 B=1+2\eta^2\E_bh^2+O(\eta^3).
\]
上下界\eqref{eq:pmdmean-kappa-bounds}于是夹出
\[
 \kappa=\frac{\eta^2}{2}\Var_b(h)+O(\eta^3).
\]
用 $W(z)=z+O(z^2)$，得到
\begin{align*}
 u_{\mathrm{mean}}(a)
 &=\eta h(a)-\kappa e^{\eta h(a)}+O(\kappa^2)\\
 &=\eta h(a)-\frac{\eta^2}{2}\Var_b(h)+O(\eta^3).
\end{align*}
精确 target 也有 $u_{\mathrm{part}}=\eta h-\log A=\eta h-\eta^2\Var_b(h)/2+O(\eta^3)$。所以小步长下两者到二阶一致，差异从更高阶开始；有限 $\eta$ 时 Lambert-$W$ 的 action-dependent 校正不同于常数 $\log A$。

在 OPD 中取 $r=\log(T/b)$，则 $h=\log(T/b)+\KL(b\Vert T)$。精确中心化 residual 的 target 是 $b^{1-\eta}T^\eta/Z$；PMD-mean 回归的 target 是式\eqref{eq:pmdmean-lambert}，仍有额外 $\chi^2$ 正则化。即使 $\eta=1$，后者也一般不是 teacher；前者则精确为 teacher。因此“小 $\eta$ 下相似”与“使用同一个算法”有不同含义。

均值 target 与最优 log-ratio 的差异属于总体目标近似。为区分这种差异与求解误差，定义
\[u^*=\log(q_{\mathrm{mean}}/b).\]
对任何正分布 $q$，平方展开和驻点条件给出精确恒等式
\begin{align}
 L_{\mathrm{mean}}(q)-L_{\mathrm{mean}}(q_{\mathrm{mean}})
 &=\tfrac12\E_b(u-u^*)^2+\E_b[(u^*-\eta h)(u-u^*)]\\
 &=\tfrac12\E_b(u-u^*)^2-\kappa\E_{q_{\mathrm{mean}}}(u-u^*)\\
 &=\tfrac12\E_b(u-u^*)^2+\kappa\KL(q_{\mathrm{mean}}\Vert q).
 \label{eq:pmdmean-excess-identity}
\end{align}
因此总体 excess loss 趋于零会使最优策略的 log-ratio 回归误差趋零。有限样本 ERM 在一致性条件下可恢复 PMD-mean 的总体 target；它仍与精确 KL target 不同。原论文 refined analysis 正是利用这类恒等式，将统计误差与目标近似造成的算法差异分开。

\section{两种方法的收益与组合条件}

TOP-D 的算术 target 提供原始 on-policy token 梯度的全概率空间二阶矩界。冻结 old policy 与 batch 允许多 epoch 数据复用，PPO ratio 机制控制某些采样方向的内部更新。后两项的全局保证仍需真实优势、采样误差和实际策略漂移控制。

MirrorOPD 的 KL 子问题提供唯一 Gibbs target、精确 projective 收缩和可展开的 inner 误差递推。它也能在 outer 内冻结行为策略并复用数据。几何收缩不会自动消除 log-ratio 负尾，算术 reward 的方差定理也不会自动推出几何收缩。

若将 TOP-D 平滑 reward 放入 KL mirror 子问题，target 会变为
\[
 q^+(a\mid s)\propto\pi_k(a\mid s)
 \left(1-\alpha+\alpha\frac{T(a\mid s)}{\pi_k(a\mid s)}\right)^\eta.
\]
$\eta=1$ 时它正好是算术混合；$0<\eta<1$ 时是旧策略到算术近端 teacher 的几何插值。该算法具有 $\alpha,\eta$ 双参数动力学，需对新算子重新证明收敛，并对实际估计量证明方差。将两套已有定理并排写出不能替代这一证明。
